\documentclass[10pt]{amsart}
\usepackage[a4paper,margin=22mm]{geometry}
\usepackage[no-math]{fontspec}
\setmainfont{Latin Modern Roman}[SmallCapsFont={lmromancaps10-regular.otf}]
\usepackage{amsmath,amssymb,amsthm,mathtools,hyperref,graphicx,xfrac,etoolbox}
\hypersetup{hidelinks}
\usepackage[scr=rsfs,cal=euler]{mathalfa}
\newcommand{\ie}{i.e.}
\newcommand{\eg}{e.g.}
\newcommand{\cf}{cf.}
\newcommand{\resp}{resp.}
\newcommand{\skpt}{\par\smallskip\noindent}
\let\Subsection\subsection
\let\Subsubsection\subsubsection
\emergencystretch=3em
\newcommand{\NoChangel}[1]{%
\expandafter\let\csname old\string#1\endcsname=#1
\let#1=\relax
\newcommand{#1}{\mathop{\mathcode`l="716C\csname old\string#1\endcsname\mathcode`l="7160}\displaylimits}%
}

\NoChangel{\log}
\NoChangel{\ln}
\NoChangel{\lim}
\NoChangel{\limsup}
\NoChangel{\liminf}
\NoChangel{\varlimsup}
\NoChangel{\varliminf}
\NoChangel{\varinjlim}
\NoChangel{\varprojlim}

\newcommand\from{\mathchoice{\longleftarrow}{\leftarrow}{\leftarrow}{\leftarrow}}
\newcommand{\To}[1]{\mathchoice{\xrightarrow{\textstyle\kern4pt#1\kern3pt}}{\stackrel{#1}{\longrightarrow}}{}{}}

\renewcommand{\theenumi}{\roman{enumi}}

\let\ts\textstyle

\newcommand{\bbullet}{{\scriptscriptstyle\bullet}}
\newcommand{\cbbullet}{{\raisebox{1pt}{$\bbullet$}}}

\usepackage{tikz}
\usepackage{tikz-cd}
\usetikzlibrary{shapes, shapes.geometric, calc, decorations.markings}
\usepackage{mathtools}

\usepackage[capitalize]{cleveref}
\crefdefaultlabelformat{#2\textup{#1}#3}
\creflabelformat{enumi}{#2\textup{(#1)}#3}

\makeatletter
\def\cref@thmoptarg[#1]#2#3#4{%
    \ifhmode\unskip\unskip\par\fi%
    \normalfont%
    \trivlist%
    \let\thmheadnl\relax%
    \let\thm@swap\@gobble%
    \thm@notefont{\fontseries\mddefault\upshape}%
    \thm@headpunct{.}% add period after heading
    \thm@headsep 5\p@ plus\p@ minus\p@\relax%
    \thm@space@setup%
    #2% style overrides
    \@topsep \thm@preskip               % used by thm head
    \@topsepadd \thm@postskip           % used by \@endparenv
    \def\@tempa{#3}\ifx\@empty\@tempa%
      \def\@tempa{\@oparg{\@begintheorem{#4}{}}[]}%
    \else%
      \refstepcounter[#1]{#3}%  <<< cleveref modification
      \@namedef{cref@#3@alias}{#1}% added
      \def\@tempa{\@oparg{\@begintheorem{#4}{\csname the#3\endcsname}}[]}%
    \fi%
    \@tempa}%
\makeatother

\newtheorem{mainthm}{Theorem}
\newtheorem{maincor}[mainthm]{Corollary}
\newtheorem{theorem}{Theorem}[section]
\newtheorem{corollary}[theorem]{Corollary}
\newtheorem{lemma}[theorem]{Lemma}
\newtheorem{proposition}[theorem]{Proposition}

\theoremstyle{definition}
\newtheorem{definition}[theorem]{Definition}
\newtheorem{example}[theorem]{Example}
\newtheorem{remark}[theorem]{Remark}

\renewcommand\themainthm{\Alph{mainthm}}

\crefname{mainthm}{Theorem}{Theorems}
\crefname{maincor}{Corollary}{Corollaries}
\crefname{theorem}{Theorem}{Theorems}
\crefname{lemma}{Lemma}{Lemmas}
\crefname{remark}{Remark}{Remarks}
\crefname{proposition}{Proposition}{Propositions}
\crefname{definition}{Definition}{Definitions}
\crefname{corollary}{Corollary}{Corollaries}
\crefname{figure}{Figure}{Figures}
\newcommand{\crefrangeconjunction}{--}
\crefname{enumi}{}{}

\let\phi\varphi
\let\eps\varepsilon
\newcommand{\lb}{\left(}
\newcommand{\rb}{\right)}
\newcommand{\id}{\operatorname{id}}

\newcommand{\CC}{\mathbb{C}}
\newcommand{\RR}{\mathbb{R}}
\newcommand{\ZZ}{\mathbb{Z}}
\newcommand{\QQ}{\mathbb{Q}}
\newcommand{\PP}{\mathbb{P}}
\newcommand{\LL}{L^\flat}
\newcommand{\hatS}{\widehat{S}}
\newcommand{\hatCO}{\widehat{\CO}{}^0_\bL}
\newcommand{\hatB}{\widehat{\calB}}
\newcommand{\hatE}{\widehat{\calE}}
\newcommand{\hatPhi}{\widehat{\Phi}}
\newcommand{\hatTheta}{\widehat{\Theta}_{\LL}}

\newcommand{\calA}{\mathcal{A}}
\newcommand{\calB}{\mathcal{B}}
\newcommand{\calC}{\mathcal{C}}
\newcommand{\calD}{\mathcal{D}}
\newcommand{\calE}{\mathcal{E}}
\newcommand{\calG}{\mathcal{G}}
\newcommand{\calL}{\mathcal{L}}
\newcommand{\calQ}{\mathcal{Q}}
\newcommand{\calT}{\mathcal{T}}
\newcommand{\Fuk}{\mathcal{F}}

\newcommand{\kk}{{\boldsymbol{k}}}
\newcommand{\bL}{{\boldsymbol{L}}}
\newcommand{\m}{\mathfrak{m}}
\newcommand{\op}{\mathrm{op}}
\newcommand{\mf}{\mathrm{mf}}
\newcommand{\diff}{\mathrm{d}}
\newcommand{\isol}{\mathrm{isol}}

\newcommand{\Char}{\operatorname{char}}
\newcommand{\rH}{\operatorname{H}}
\newcommand{\HH}{\operatorname{HH}}
\newcommand{\rCC}{\operatorname{CC}}
\newcommand{\HF}{\operatorname{HF}}
\newcommand{\CF}{\operatorname{CF}}
\newcommand{\Spec}{\operatorname{Spec}}
\newcommand{\LMF}[1][]{\operatorname{\mathcal{LM}}_{#1}}
\newcommand{\End}{\operatorname{End}}
\newcommand{\CO}{\operatorname{\mathcal{CO}}}
\newcommand{\OC}{\operatorname{\mathcal{OC}}}
\newcommand{\QH}{\operatorname{QH}}
\newcommand{\SH}{\operatorname{SH}}
\newcommand{\Jac}{\operatorname{Jac}}
\newcommand{\Jacisol}{\operatorname{Jac}_\isol}
\newcommand{\COisol}{\CO^0_{\bL, \isol}}
\newcommand{\Tw}{\operatorname{Tw}}
\newcommand{\oPi}{\operatorname{\Pi}}
\newcommand{\PiTw}{\oPi \Tw}

\newcommand{\qprod}{\mathbin{\star}}
\newcommand{\Shk}[2]{\langle #1, #2 \rangle_\mathrm{Shk}}
\newcommand{\Poinc}[2]{\langle #1, #2 \rangle_X}


\let\sourceLocalCref\cref
\renewcommand{\cref}[1]{\ifstrequal{#1}{TheoremA}{\href{https://doi.org/10.5802/jep.290}{Theorem~A (source)}}{\ifstrequal{#1}{CorollaryB}{\href{https://doi.org/10.5802/jep.290}{Corollary~B (source)}}{\sourceLocalCref{#1}}}}
\title[Generation from smoothness]{Automatic split-generation and the quantum-cohomology factor from a smooth Lagrangian endomorphism algebra}
\author{Jack Smith}
\date{Published version; DOI 10.5802/jep.290}
\begin{document}
\maketitle
\section*{Source conventions for the Fukaya category}
\setcounter{section}{1}\setcounter{theorem}{7}
\begin{remark}
\label{rmkFukTechnical}
\begin{enumerate}
\item Each Lagrangian in the Fukaya category should be equipped with a $\ZZ/2$-grading with respect to a fixed $\ZZ/2$-grading of $X$, in~the sense of \cite{SeidelGraded}. If~we only care about consequences for $\QH^*(X)$ then we can equip $X$ with its canonical $\ZZ/2$-grading, in~which case a $\ZZ/2$-grading of a Lagrangian is just an orientation (in particular, the torus $L$ has a $\ZZ/2$-grading).
\item Each Lagrangian should also be equipped with a relative spin structure, to~determine the signs of pseudoholomorphic disc counts. In~particular, the superpotential $W_L$ of the torus $L$ depends on the choice of relative spin structure on~$L$. Happily there is a canonical choice in this case, namely the unique spin structure induced by picking an arbitrary identification $L \!\cong\! \RR^n / \ZZ^n$ and using this to trivialise~$TL$.
\item To do monotone Floer theory with multiple Lagrangians, we need to assume that the fundamental group of each has trivial image in $\pi_1(X)$. For all statements involving the Fukaya category, therefore, we should impose the additional hypothesis on each Lagrangian $K$ that $\pi_1(K)$ has trivial image in $\pi_1(X)$. (Alternatively one could weaken these conditions to be about $\rH_1$ instead of $\pi_1$, at the expense of strengthening the monotonicity condition to be in terms of relative $\rH_2$ rather than relative $\pi_2$.) \cref{TheoremA}, and therefore also \cref{CorollaryB}, does not require this hypothesis since it only relies on Floer theory of $L$ with itself.
\item We can twist $\QH^*(X)$ and the Fukaya category by any homomorphism
\[
B : \rH_2(X; \ZZ) \to \kk^\times,
\]
and our results have an analogue where the local system on each Lagrangian $K$ is replaced by a lift of $B$ to $\rH_2(X, K; \ZZ) \to \kk^\times$. We~need to assume that such a lift exists for our torus $L$, \ie that that restriction of $B$ to the image of $\rH_2(L; \ZZ)$ is trivial. We~leave the details to the interested reader.
\end{enumerate}
\end{remark}
\setcounter{section}{1}\setcounter{equation}{1}

\section{Generation from smoothness}
\label{secGeneration}

\subsection{Split-generation}
\label{sscGeneration}

First we recap various categorical preliminaries, closely following \cite{GanatraAutomaticGeneration}.

Recall that an $A_\infty$-category $\calC$ has a \emph{pretriangulated envelope}, denoted $\Tw \calC$, and an \emph{idempotent completion} or \emph{split-closure}, denoted $\oPi \calC$. Each of these is well-defined up to quasi-equivalence and carries a cohomologically full and faithful embedding of~$\calC$. An explicit model for $\Tw \calC$ is the category of \emph{twisted complexes} (hence the name), given by embedding $\calC$ in $\calC\text{-mod}$ using Yoneda, and then taking the full subcategory obtained from the image by iteratively taking mapping cones.

We will mostly be interested in the split-closed pretriangulated envelope, $\PiTw \calC$, sometimes also denoted $\mathrm{perf}(\calC)$ (\eg in \cite{GanatraAutomaticGeneration}). For example, $\Fuk(X)_\lambda$ implicitly means the split-closed pretriangulated envelope of the Fukaya category whose objects are actual Lagrangians. The embedding
\[
\PiTw \calC \to \PiTw(\PiTw\calC)
\]
is a quasi-equivalence, and we say that a categorical property is \emph{Morita-invariant} if it is preserved by quasi-equivalences and by applying $\PiTw$. Similarly, a construction is Morita-invariant if quasi-equivalences and embeddings $\calC \to \PiTw \calC$ induce (quasi\nobreakdash-)isomorphisms on it. For example, Hochschild homology is Morita-invariant since a functor $F : \calC \to \calD$ induces a homomorphism $\HH_*(\calC) \to \HH_*(\calD)$ which is an isomorphism if $F$ is a quasi-equivalence or an embedding of the form $\calC \to \PiTw \calC$. Similarly, Hochschild cohomology is Morita-invariant, in~the sense that a functor $F : \calC \to \calD$ induces a zigzag
\begin{equation}
\label{eqHHmaps}
\HH^*(\calC) \to \HH^*(\calC, \calD) \from \HH^*(\calD)
\end{equation}
and both of these arrows are isomorphisms if $F$ is a quasi-equivalence or an embedding $\calC \to \PiTw \calC$.

\begin{remark}
We refer the reader to \cite[\S 2]{RitterSmith} and \cite[App.\,A]{SheridanFano} for the fundamentals of Hochschild homology and cohomology in the setting of Fukaya categories. In~\eqref{eqHHmaps} the group $\HH^*(\calC, \calD)$ denotes Hochschild cohomology of $\calC$ with coefficients in $\calD$, viewed as a $\calC$-bimodule via $F$, and the groups $\HH^*(\calC)$ and $\HH^*(\calD)$ are really shorthand for $\HH^*(\calC, \calC)$ and $\HH^*(\calD, \calD)$. The first arrow is pushforward on coefficients, whilst the second arrow is pullback on Hochschild cochains.
\end{remark}

If $\calD$ is a full subcategory of $\calC$ then there is a cohomologically full and faithful embedding of $\PiTw \calD$ into $\PiTw \calC$. We~say that $\calD$ \emph{split-generates} $\calC$ if this embedding is a quasi-equivalence.

\subsection{The Shklyarov pairing}
\label{sscShklyarov}

Recall that a dg- or $A_\infty$-category over $\kk$ is \emph{proper} if the total cohomology of each morphism space is finite-dimensional. Any such category has a \emph{Shklyarov pairing} on its Hochschild homology, denoted $\Shk{{-}}{{-}}$. This was introduced by Shklyarov \cite{ShklyarovHirzebruch} for dg-algebras (\ie dg-categories with a single object) and extended to $A_\infty$-categories by Sheridan \cite[\S 5]{SheridanHodge}, who calls it the \emph{Mukai} pairing. Sheridan proves that this pairing is Morita-invariant \cite[Prop.\,5.20]{SheridanHodge}, in~the sense that it is preserved by the isomorphism on Hochschild homology induced by a quasi-equivalence or an embedding $\calC \to \PiTw \calC$. It~follows immediately from his explicit formula for the pairing \cite[(35)]{SheridanHodge} that the pairing is also preserved under inclusions of full subcategories, and that if $\calA_1, \calA_2 \subset \calC$ are full subcategories which are categorically orthogonal (morphism complexes between $\calA_1$ and $\calA_2$ are acyclic) then the images of $\HH_*(\calA_1)$ and $\HH_*(\calA_2)$ in $\HH_*(\calC)$ are orthogonal with respect to it.

Recall next that a dg- or $A_\infty$-category over $\kk$ is \emph{(homologically) smooth} if its diagonal bimodule is split-generated by Yoneda bimodules. Shklyarov shows that if a dg-algebra is smooth and proper then its Shklyarov pairing is non-degenerate \cite[Th.\,5.3]{ShklyarovHirzebruch}, and that its Hochschild homology is finite-dimensional \cite[Th.\,4.6]{ShklyarovSerre}, so the pairing is in fact perfect. Since smoothness is Morita-invariant \cite[Prop.\,20]{GanatraAutomaticGeneration}, and every $A_\infty$-algebra is quasi-isomorphic to a dg-algebra, it~follows from Sheridan's results that the same properties hold for smooth and proper $A_\infty$-algebras.

\begin{proposition}[{\cite[Prop.\,5.24]{SheridanHodge}}]
\label{propSmoothProper}
If an $A_\infty$-algebra $\calA$ is smooth and proper then $\HH_*(\calA)$ is finite-dimensional and its Shklyarov pairing is perfect.\qed
\end{proposition}

\subsection{The closed--open and open--closed string maps}
\label{sscStringMaps}

Now we focus on the case of interest, where $\calC$ is the Fukaya category $\Fuk(X)_\lambda$. Recall from \cite[\S\S 2.5--2.6]{SheridanFano} that there is \emph{closed--open string map}
\[
\CO_\lambda : \QH^*(X) \to \HH^*(\Fuk(X)_\lambda).
\]
which is a $\ZZ/2$-graded unital $\kk$-algebra homomorphism. There is likewise an \emph{open--closed string map}
\[
\OC_\lambda : \HH_*(\Fuk(X)_\lambda) \to \QH^{*+n}(X),
\]
which is a $\ZZ/2$-graded $\QH^*(X)$-module map, where $\alpha \in \QH^*(X)$ is defined to act on $\HH_*(X)$ via the cap product action of $\CO_\lambda(\alpha)$. Recalling that $c_1 \qprod$ denotes the quantum multiplication action of $c_1(X)$ on $\QH^*(X)$, we have the following.

\begin{proposition}[{\cite[Th.\,9.5(1)]{RitterSmith}, or~\cite[Cor.\,2.11]{SheridanFano} over $\CC$}]
\label{propEspace}
Assume that $\Char \kk \neq 2$ or that we only allow orientable Lagrangians in the Fukaya category. Then the image of $\OC_\lambda$ is contained in the generalised $\lambda$-eigenspace $\QH^*(X)_\lambda$ of $c_1 \qprod$.\qed
\end{proposition}

\begin{remark}
There is a corresponding result for $\CO_\lambda$, but we will not use it. Over $\CC$ this states that $\CO_\lambda$ vanishes on $\QH^*(X)_\mu$ if $\mu \neq \lambda$; see \cite[Prop.\,2.9]{SheridanFano} or \cite[Th.\,9.6]{RitterSmith}. In~general (\ie without assuming that $\kk$ is algebraically closed) it states that $\CO_\lambda$ vanishes on the unique $c_1 \qprod$-invariant complement $C$ to $\QH^*(X)_\lambda$ in $\QH^*(X)$. This can be proved by factorising the characteristic polynomial $\chi(T)$ of $c_1\qprod$ as $(T-\lambda)^m p(T)$ for some non-negative integer $m$ and some polynomial $p$ with $p(\lambda) \neq 0$. By~Bézout's lemma there exist polynomials $f$ and $g$ such that $1 = (T-\lambda)^mf(T) + p(T)g(T)$. The maps $(c_1 - \lambda)^mf(c_1) \qprod$ and $p(c_1)g(c_1)\qprod$ then represent projection onto $C$ and $\QH^*(X)_\lambda$ respectively. One argues, as in \cite[Prop.\,2.9]{SheridanFano}, that $\CO_{\lambda}(p(c_1))$ is invertible in $\HH^*(\Fuk(X)_\lambda)$. On the other hand, by Cayley--Hamilton we have $\chi(c_1) = 0$ in $\QH^*(X)$, and hence
\begin{equation}
\label{eqCOespace}
0 = \CO_\lambda(\chi(c_1))= \CO_\lambda((c_1 - \lambda)^m p(c_1)) = \CO_\lambda((c_1 - \lambda)^m) \cup \CO_\lambda(p(c_1)),
\end{equation}
where $\cup$ is the cup product on $\HH^*(\Fuk(X)_\lambda)$. Combining \eqref{eqCOespace} with the fact that $\CO_\lambda(p(c_1))$ is invertible, we deduce that $\CO_\lambda((c_1 - \lambda)^m) = 0$. Therefore $\CO_\lambda$ annihilates all multiples of $(c_1-\lambda)^mf(c_1)$, and hence annihilates $C$.
\end{remark}

The category $\Fuk(X)_\lambda$ is proper, so its Hochschild homology comes equipped with its Shklyarov pairing. Meanwhile, quantum cohomology carries the Poincaré pairing $\Poinc{{-}}{{-}}$. Crucially, $\OC_\lambda$ is compatible with these pairings in the following sense.

\begin{proposition}[{\cite[Th.\,31, credited to Ganatra--Perutz--Sheridan]{GanatraAutomaticGeneration}}]
\label{propOCisometry}
For any classes $a$ and $b$ in $\HH_*(\Fuk(X)_\lambda)$ we have
\begin{flalign*}
&& \Shk{a}{b} &= (-1)^{n(n+1)/2}\Poinc{\OC_\lambda(a)}{\OC_\lambda(b)}. && \qed
\end{flalign*}
\end{proposition}

Given a full subcategory $\calA \subset \Fuk(X)_\lambda$ we can compose $\CO_\lambda$ with the restriction map
\[
\HH^*(\Fuk(X)_\lambda) \to \HH^*(\calA)
\]
to give a homomorphism $\CO_\calA : \QH^*(X) \to \HH^*(\calA)$. Similarly we can compose $\OC_\lambda$ with the homomorphism $\HH_*(\calA) \to \HH_*(\Fuk(X)_\lambda)$ induced by the inclusion functor to give a homomorphism $\OC_\calA : \HH_*(\calA) \to \QH^{*+n}(X)$.

\subsection{The weak proper Calabi--Yau structure}
\label{sscwpCY}

Sheridan shows \cite[\S 2.8]{SheridanFano} that $\Fuk(X)_\lambda$ also carries an \emph{$n$-dimensional weak proper Calabi--Yau (wpCY) structure} $[\phi] \in \HH_n(\Fuk(X)_\lambda)^\vee$.

\begin{lemma}[{\cite[Lem.\,A.2]{SheridanFano}}]
\label{lemwpCY}
This induces a perfect pairing
\[
\langle{-}, {-}\rangle_\mathrm{wpCY} : \HH^*(\Fuk(X)_\lambda) \times \HH_{n-*}(\Fuk(X)_\lambda) \to \kk \quad \text{via} \quad \langle \psi, b \rangle_\mathrm{wpCY} = [\phi] (\psi \cap b).
\]
The same arguments show that the pairing remains perfect after restricting to a full subcategory $\calA$ on a subset of Lagrangians. We~denote this restricted pairing by $\langle{-}, {-}\rangle_{\mathrm{wpCY}, \calA}$.\qed
\end{lemma}

He then shows that $\CO_\lambda$ and $\OC_\lambda$ are adjoint with respect to this pairing, in~the following sense.

\begin{proposition}[{\cite[Prop.\,2.6]{SheridanFano}}]
\label{propCOOCduality}
For $\alpha \!\in\! \QH^*(X)$ and $b \!\in\! \HH_{n-*}(\Fuk(X)_\lambda)$ we~have
\[
\langle \CO_\lambda(\alpha), b\rangle_\mathrm{wpCY} = \langle \alpha, \OC_\lambda(b) \rangle_X.
\]
By the same arguments, for any full subcategory $\calA \subset \Fuk(X)_\lambda$ on a subset of Lagrangians, the maps $\CO_\calA$ and $\OC_\calA$ are likewise adjoint with respect to $\langle {-}, {-}\rangle_{\mathrm{wpCY}, \calA}$.\qed
\end{proposition}

\subsection{Decompositions and the generation}

For each object $T$ in $\Fuk(X)_\lambda$ let
\[
\CO_T : \QH^*(X) \to \HH^*(\hom^*(T, T))
\]
denote the restriction of $\CO_\lambda$ to $T$; this is just $\CO_\calA$ in the case where $\calA$ contains a single object, $T$. Let
\[
\CO^0_T : \QH^*(X) \to \rH^*(\hom^*(T, T))
\]
denote the projection of $\CO_T$ to length-zero Hochschild cochains. Given an idempotent $e \in \QH^*(X)$ of even degree, or~equivalently a factor $Q = e \qprod \QH^*(X)$ of $\QH^*(X)$, there is an associated full subcategory $\Fuk(X)_{\lambda,Q}$ of $\Fuk(X)_\lambda$, comprising those objects~$T$ for which $\CO^0_T(e) = 1_T \in \rH^*(\hom^*(T, T))$.

\begin{theorem}
\label{thmGen}
Suppose $\LL$ is a Lagrangian in $\Fuk(X)_\lambda$, and that its endomorphism $A_\infty$-algebra $\calA$ is smooth. The following then hold:
\begin{enumerate}
\item\label{genOC} $\OC_\calA$ gives an isomorphism from $\HH_*(\calA)$ onto its image $Q \subset \QH^*(X)$.
\item\label{genEspace} If $\Char \kk \neq 2$ or $L$ is orientable then $Q$ is contained in the generalised $\lambda$\nobreakdash-eigen\-space, $\QH^*(X)_\lambda$, of $c_1\qprod$.
\item\label{genVSdecomp} As a $\ZZ/2$-graded vector space, $\QH^*(X)$ is the internal direct sum $Q \oplus Q^\perp$.
\item\label{genAlgDecomp} Both $Q$ and $Q^\perp$ are two-sided ideals in $\QH^*(X)$, so we obtain a decomposition of $\ZZ/2$-graded algebras $\QH^*(X) = Q \times Q^\perp$.
\item\label{genIdemp} The $Q$-component $e$ of $1_X \in \QH^*(X)$ is an even-degree idempotent, satisfying $Q = e \qprod \QH^*(X)$ and $Q^\perp = (1_X-e) \qprod \QH^*(X)$.
\item\label{genCO} $\CO_{\LL} \!=\! \CO_\calA$ induces an algebra isomorphism $Q \!\to\! \HH^*(\calA)$, and \hbox{$\ker \CO_{\LL} \!=\! Q^\perp$}.
\item\label{genSplitGen} $\LL$ lies in $\Fuk(X)_{\lambda, Q}$ and split-generates it.
\end{enumerate}
\end{theorem}

\skpt
\begin{proof}
\eqref{genOC} The algebra $\calA$ is smooth by assumption, and proper because we're dealing with compact Lagrangians, so by \cref{propSmoothProper} $\HH_*(\calA)$ is finite-dimensional and its Shklyarov pairing is perfect. Since $\OC_\lambda$ is an isometry with respect to this pairing, in~the sense of \cref{propOCisometry}, we deduce that it is injective on $\HH_*(\calA)$ and hence induces an isomorphism onto $Q$.

\eqref{genEspace} This follows immediately from \cref{propEspace}.

\eqref{genVSdecomp} From the proof of \eqref{genOC}, the restriction of the Poincaré pairing to $Q$ is perfect. The vector space decomposition of $\QH^*(X)$ then follows from the elementary linear algebra result that if $V$ is a finite-dimensional vector space, $\beta$ is an arbitrary bilinear form on $V$, and $W \subset V$ is a subspace on which $\beta$ is non-degenerate, then $V = W \oplus W^{\perp_\beta}$. The fact that $Q$ is actually a $\ZZ/2$-graded subspace is an immediate consequence of $\OC_\lambda$ being a $\ZZ/2$-graded map. $Q^\perp$ is then automatically $\ZZ/2$-graded since the Poincaré pairing is $\ZZ/2$-graded. Note also that since the Poincaré pairing is symmetric on even-degree elements and skew-symmetric on odd-degree elements we have $Q^\perp = {}^\perp Q$.

\eqref{genAlgDecomp} Suppose we can show that $Q$ is a left ideal. The Frobenius algebra property of $\QH^*(X)$ then tells us that ${}^\perp Q$ ($=Q^\perp$) is a right ideal. Since both $Q$ and $Q^\perp$ are $\ZZ/2$-graded subspaces, graded-commutativity of $\QH^*(X)$ then implies that both are in fact two-sided ideals, and we're done.

It therefore suffices to show that $Q$ is a left ideal, so take $q \in Q$ and $\alpha \in \QH^*(X)$. We~wish to show that $\alpha q \in Q$. By~definition of $Q$ there exists $a \in \HH_*(\calA)$ with $q = \OC_\calA(a)$. We~can write this as $\OC_\lambda(i_*a)$, where $i_* : \HH_*(\calA) \to \HH_*(\Fuk(X)_\lambda)$ is the homomorphism induced by inclusion $\calA \to \Fuk(X)_\lambda$. Using the fact that $\OC_\lambda$ is a $\QH^*(X)$-module map we obtain
\[
\alpha q = \alpha \OC_\lambda(i_*a) = \OC_\lambda(\CO_\lambda(\alpha) \cap i_*a).
\]
The image of $i_*$ is manifestly a $\HH^*(\Fuk(X)_\lambda)$-submodule of $\HH_*(\Fuk(X)_\lambda)$, so it contains $\CO_\lambda(\alpha) \cap i_*a$. This means that $\alpha q$ is in the image of $\OC_\lambda \circ i_* = \OC_\calA$, which is pre\-cisely~$Q$, as wanted.

\eqref{genIdemp} It follows from \eqref{genVSdecomp} that $1_X$ splits into even-degree pieces $e \!\in\! Q$ and \hbox{$1_X\!-\!e\!\in\! Q^\perp$}. The algebra decomposition in \eqref{genAlgDecomp} then tells us that $e (1_X - e) = 0$, so $e$ is idempotent. It~also tells us that $e$ and $1_X - e$ are the units in $Q$ and $Q^\perp$ respectively, so $Q = e \QH^*(X)$ and $Q^\perp = (1-e) \QH^*(X)$.

\eqref{genCO} Let $\eps \in \HH_*(\calA)$ be the unique element satisfying $e = \OC_\calA(\eps)$. Again writing $\OC_\calA$ as $\OC_\lambda \circ i_*$ and using the fact that $\OC_\lambda$ is a $\QH^*(X)$-module map, we see that for all $\alpha \in Q$
\[
\alpha = \alpha e = \OC_\lambda(\CO_\lambda(\alpha) \cap i_*\eps) = \OC_\lambda(i_*(\CO_\calA(\alpha) \cap \eps)).
\]
This forces $\CO_{\LL} = \CO_\calA$ to be injective on $Q$. By~now combining \eqref{genOC} with the equality $\dim \HH^*(\calA) = \dim \HH_{n-*}(\calA)$ resulting from \cref{lemwpCY} we obtain $\dim \HH^*(\calA) = \dim Q$, and deduce that $\CO_{\LL}$ is actually an isomorphism $Q \to \HH^*(\calA)$.

It remains to show that $\ker \CO_{\LL} = Q^\perp$, and since we know
\[
\dim Q^\perp = \dim \QH^*(X) - \dim Q \leq \dim \QH^*(X) - \operatorname{rank} \CO_{\LL} = \dim \ker \CO_{\LL}
\]
it is enough to show that $\ker \CO_{\LL} \subset Q^\perp$. Suppose then that $\alpha \in \ker \CO_{\LL}$. Since $\CO_{\LL}$ is an algebra homomorphism we get $\CO_{\LL}(e \alpha) = 0$. But $e \alpha$ is in $Q$, so by injectivity of $\CO_{\LL}$ on $Q$ we deduce that $e \alpha = 0$, and hence $\alpha = (1_X - e) \alpha$. Thus $\alpha \in (1_X - e) \QH^*(X) = Q^\perp$, and so $\ker \CO_{\LL} \subset Q^\perp$.

\eqref{genSplitGen} To prove that $\LL$ lies in $\Fuk(X)_{\lambda,Q}$ we need to show that $\CO^0_{\LL}(e) = 1_{\LL}$, for which it suffices to show that $\CO_{\LL}(e)$ is the unit $1_{\HH^*(\calA)}$ in $\HH^*(\calA)$. To do this, note that since $\CO_{\LL}$ induces an isomorphism $Q \to \HH^*(\calA)$ there exists a unique $e' \in Q$ such that $\CO_{\LL}(e') = 1_{\HH^*(\calA)}$. We~then have
\[
\CO_{\LL}(e') = \CO_{\LL}(e e') = \CO_{\LL}(e) \cup \CO_{\LL}(e') = \CO_{\LL}(e).
\]
Injectivity of $\CO_{\LL}$ on $Q$ gives $e' = e$, and hence $\CO_{\LL}(e) = 1_{\HH^*(\calA)}$, as wanted.

For any other object $T$ in $\Fuk(X)_{\lambda, Q}$, we have $\CO^0_T(e) = 1_T$ by definition. So~$\CO^0_T \circ \OC_\calA(\eps) = 1_T$, where $\eps \in \HH_*(\calA)$ is the unique element with $\OC_\calA(\eps) = e$, as above. The fact that $T$ is split-generated by $\LL$ then follows from Abouzaid's generation criterion \cite[Prop.\,1.3, Lem.\,1.4]{AbouzaidGeometricCriterion}.
\end{proof}
\begin{thebibliography}{99}
\bibitem[1]{AbouzaidGeometricCriterion} M. Abouzaid - “A geometric criterion for generating the Fukaya category”, Publ. Math. Inst. Hautes Études Sci. (2010) no. 112, p. 191-240.
\bibitem[15]{GanatraAutomaticGeneration} S. Ganatra - “Automatically generating Fukaya categories and computing quantum cohomology”, 2016, arXiv:1605.07702.
\bibitem[21]{RitterSmith} A. F. Ritter \& I. Smith - “The monotone wrapped Fukaya category and the open-closed string map”, Selecta Math. (N.S.) 23 (2017) no. 1, p. 533-642.
\bibitem[23]{SeidelGraded} P. Seidel - “Graded Lagrangian submanifolds”, Bull. Soc. math. France 128 (2000) no. 1, p. 103-149.
\bibitem[25]{SheridanFano} N. Sheridan - “On the Fukaya category of a Fano hypersurface in projective space”, Publ. Math. Inst. Hautes Études Sci. 124 (2016), p. 165-317.
\bibitem[26]{SheridanHodge} N. Sheridan - “Formulae in noncommutative Hodge theory”, J. Homotopy Relat. Struct. 15 (2020) no. 1, p. 249-299.
\bibitem[27]{ShklyarovSerre} D. Shklyarov - “On Serre duality for compact homologically smooth DG algebras”, 2007, arXiv:math/0702590.
\bibitem[28]{ShklyarovHirzebruch} D. Shklyarov - “Hirzebruch-Riemann-Roch-type formula for DG algebras”, Proc. London Math. Soc. (3) 106 (2013) no. 1, p. 1-32.
\end{thebibliography}
\end{document}
